% Luc Destombe --- licence CC BY-NC-SA 4.0
% https://creativecommons.org/licenses/by-nc-sa/4.0/deed.fr
% Fichier autonome : compiler avec pdflatex, deux fois (signets, nombre de pages).
\documentclass[11pt,a4paper]{article}
\makeatletter
\newif\ifeleve
\elevefalse

\newif\iflycee@palatino
\lycee@palatinotrue

\RequirePackage[utf8]{inputenc}
\RequirePackage[T1]{fontenc}
\RequirePackage[french]{babel}
\RequirePackage{etoolbox}

\newcommand{\ShorthandsOff}{\@ifundefined{shorthandoff}{}{\shorthandoff{;:!?}}}
\newcommand{\ShorthandsOn}{\@ifundefined{shorthandon}{}{\shorthandon{;:!?}}}
\ShorthandsOff
\AtBeginDocument{\ShorthandsOn}
\AtBeginEnvironment{tikzpicture}{\ShorthandsOff}

\RequirePackage{geometry}
\RequirePackage{fancyhdr}
\RequirePackage{lastpage}
\RequirePackage{multicol}
\RequirePackage{array}
\RequirePackage{enumitem}
\RequirePackage{xcolor}
\RequirePackage{needspace}

\RequirePackage{amsmath,amssymb}
\RequirePackage{mathpazo}
\DeclareMathSizes{10.95}{10.95}{8}{5.5}
\begingroup\catcode`\_=\active \catcode`\^=\active
\gdef_#1{\sb{\scriptscriptstyle #1}}
\gdef^#1{\sp{\scriptscriptstyle\lycee@moinsepais #1}}
\endgroup
\newcommand\lycee@moins{\mathbin{%
  \setbox\z@\hbox{$\m@th\scriptscriptstyle\lycee@moinsorig$}%
  \dimen@=.55\wd\z@ \advance\dimen@-.5pt
  \hbox to.75\wd\z@{\hss$\m@th\scriptscriptstyle\vcenter{\hbox{%
    \tikz\draw[line width=.5pt,line cap=round](0,0)--(\the\dimen@,0);}}$\hss}}}
\begingroup\lccode`\~=`\- \lowercase{\endgroup\let~}\lycee@moins
\newcommand\lycee@moinsepais{\mathcode`\-="8000 }
\AtBeginDocument{\mathchardef\lycee@moinsorig=\mathcode`\-\relax}
\AtBeginDocument{\catcode`\_=12 \mathcode`\_="8000
  \catcode`\^=12 \mathcode`\^="8000 }
\newcommand\lycee@exposants{%
  \fontdimen13\textfont2=.48\fontdimen6\textfont2
  \fontdimen14\textfont2=.48\fontdimen6\textfont2
  \fontdimen15\textfont2=.48\fontdimen6\textfont2
  \fontdimen13\scriptfont2=.48\fontdimen6\scriptfont2
  \fontdimen14\scriptfont2=.48\fontdimen6\scriptfont2
  \fontdimen15\scriptfont2=.48\fontdimen6\scriptfont2}
\every@math@size\expandafter{\the\every@math@size\lycee@exposants}
\newlength{\retraitcalculs}
\setlength{\retraitcalculs}{2em}

\RequirePackage{tikz}
\usetikzlibrary{arrows.meta,calc,babel,shapes.misc}
\RequirePackage[most]{tcolorbox}
\tcbuselibrary{skins,breakable}

\definecolor{coulDef}{HTML}{1F4E79}
\definecolor{coulProp}{HTML}{7030A0}
\definecolor{coulRem}{HTML}{C55A11}
\definecolor{coulEx}{HTML}{2E7D32}
\definecolor{coulExo}{HTML}{B03A2E}
\definecolor{coulPart}{HTML}{1F4E79}
\definecolor{coulMeth}{HTML}{1B6E4F}
\definecolor{coulCourbe}{HTML}{0B5ED7}
\definecolor{coulCourbeB}{HTML}{C00000}
\definecolor{coulCourbeC}{HTML}{2E7D32}
\definecolor{coulGrille}{HTML}{8C8C8C}
\definecolor{coulRep}{HTML}{1B6E4F}
\definecolor{coulBar}{HTML}{2E7D32}

\newcommand{\R}{\mathbb{R}}
\newcommand{\Rs}{\mathbb{R}^{*}}
\renewcommand{\le}{\leqslant}
\renewcommand{\ge}{\geqslant}
\newcommand{\vect}[1]{\overrightarrow{#1}}
\newcommand{\norme}[1]{\left\lVert #1 \right\rVert}
\newcommand{\intervalle}[2]{\left[#1\,;#2\right]}
\RequirePackage{cancel}
\renewcommand{\CancelColor}{\color{coulExo}}
\newcommand{\fois}[1]{\times{\color{coulExo}#1}}
\newcommand{\barre}[1]{\cancel{#1}}

\tikzset{grille/.style={coulGrille,line width=0.4pt}}
\newcommand{\quadrillage}[4]{\draw[grille,step=1] (#1,#2) grid (#3,#4);}
\tikzset{
  croix/.style={cross out,draw,line width=0.6pt,inner sep=0pt,minimum size=4pt},
  croixdroite/.style={croix,rotate=45,line width=0.8pt},
}
\newcommand{\pt}[3]{\path (#1) node[croix]{} node[#2,font=\small,inner sep=2.5pt]{$#3$};}

\newcommand{\lycee@repere}[4]{%
  \draw[grille,step=1] (#1,#2) grid (#3,#4);
  \draw[-{Stealth[length=2mm]},thick] (#1,0)--({#3+0.4},0) node[below left=-1pt]{$x$};
  \draw[-{Stealth[length=2mm]},thick] (0,#2)--(0,{#4+0.4}) node[below left=-1pt]{$y$};
  \node[below left,font=\scriptsize,inner sep=1.5pt] at (0,0) {$0$};}
\newcommand{\lycee@gradx}[1]{\foreach \k/\l in {#1}{%
  \draw (\k,0.07)--(\k,-0.07) node[below,font=\scriptsize,inner sep=1.5pt]{$\l$};}}
\newcommand{\lycee@grady}[1]{\foreach \k/\l in {#1}{%
  \draw (0.07,\k)--(-0.07,\k) node[left,font=\scriptsize,inner sep=1.5pt]{$\l$};}}
\newcommand{\lycee@intff}[2]{\left[#1\,;#2\right]}
\newcommand{\lycee@intfo}[2]{\left[#1\,;#2\right[}
\newcommand{\lycee@intof}[2]{\left]#1\,;#2\right]}
\newcommand{\lycee@intoo}[2]{\left]#1\,;#2\right[}
\newcommand{\lycee@ens}[1]{\left\{#1\right\}}
\AtBeginDocument{%
  \providecommand{\repere}{\lycee@repere}%
  \providecommand{\gradx}{\lycee@gradx}%
  \providecommand{\grady}{\lycee@grady}%
  \providecommand{\Cf}{\mathcal{C}\sb{\scriptscriptstyle f}}%
  \providecommand{\Cg}{\mathcal{C}\sb{\scriptscriptstyle g}}%
  \providecommand{\intff}{\lycee@intff}%
  \providecommand{\intfo}{\lycee@intfo}%
  \providecommand{\intof}{\lycee@intof}%
  \providecommand{\intoo}{\lycee@intoo}%
  \providecommand{\ens}{\lycee@ens}}

\newlist{questions}{enumerate}{3}
\setlist[questions,1]{label=\arabic*\textdegree),leftmargin=*,itemsep=2pt,topsep=3pt}
\setlist[questions,2]{label=\alph*),leftmargin=*,itemsep=1pt,topsep=2pt}
\setlist[questions,3]{label=\roman*),leftmargin=*,itemsep=1pt,topsep=1pt}
\setlist[itemize]{label=\textbullet,leftmargin=*,itemsep=2pt,topsep=3pt}

\newcommand{\besoinplace}[1]{\par\penalty\@M\begingroup
  \dimen@=#1\relax\dimen@ii=\pagegoal\advance\dimen@ii-\pagetotal
  \ifdim\dimen@>\dimen@ii\ifdim\dimen@ii>\z@\vfil\fi\break\fi\endgroup}

\newcommand{\lycee@pied}{\thepage/\pageref{LastPage}}
\newcommand{\entete}[2]{%
  \pagestyle{fancy}\fancyhf{}%
  \fancyhead[L]{\footnotesize\color{coulPart}\textbf{#1}}%
  \fancyhead[R]{\footnotesize\color{coulPart}\textbf{#2}}%
  \fancyfoot[C]{\footnotesize\lycee@pied}%
  \renewcommand{\headrulewidth}{0.4pt}%
  \renewcommand{\headrule}{\hbox to\headwidth{%
    \color{coulPart!40}\leaders\hrule height \headrulewidth\hfill}}}

\RequirePackage{ccicons}
\newcommand{\lycee@licence}{{\scriptsize\color{gray}%
  \copyright\ Luc Destombe --- \ccbyncsa\ CC BY-NC-SA 4.0}}
\newif\iflycee@licence \lycee@licencetrue
\newcommand{\sanslicence}{\lycee@licencefalse}
\AtBeginDocument{\iflycee@licence\AddToHookNext{shipout/foreground}{%
  \put(0,0){\rlap{\kern\dimexpr1in+\hoffset+\oddsidemargin+\textwidth\relax
    \llap{\raisebox{-\dimexpr1in+\voffset+\topmargin+\headheight+\headsep
      +\textheight+\footskip\relax}{\lycee@licence}}}}}\fi}

\newcommand{\formulesserrees}{%
  \setlength{\abovedisplayskip}{3pt}\setlength{\belowdisplayskip}{3pt}%
  \setlength{\abovedisplayshortskip}{2pt}\setlength{\belowdisplayshortskip}{3pt}}

\setlength{\parindent}{0pt}

\RequirePackage[implicit=false,hidelinks,pdfencoding=auto,psdextra,bookmarksopen,
  bookmarksopenlevel=1,pdfcreator={},pdfproducer={}]{hyperref}
\RequirePackage{bookmark}
\hypersetup{pdfauthor={Luc Destombe},pdfkeywords={CC BY-NC-SA 4.0}}
\newcounter{lycee@signet}
\newcommand{\signet}[3][]{\stepcounter{lycee@signet}%
  \edef\lycee@dest{\if\relax\detokenize{#1}\relax
    signet.\arabic{lycee@signet}\else#1\fi}%
  \hypertarget{\lycee@dest}{}\bookmark[level=#2,dest=\lycee@dest]{#3}}
\pdfstringdefDisableCommands{%
  \def\R{R}\def\vect#1{#1}\def\Cf{Cf}\def\Cg{Cg}\def\fois#1{×#1}%
  \def\barre#1{#1}\def\intervalle#1#2{[#1;#2]}\def\intff#1#2{[#1;#2]}%
  \def\textsuperscript#1{#1}\def\dfrac#1#2{#1/#2}\def\frac#1#2{#1/#2}%
  \def\vec#1{#1⃗}}%
\geometry{top=2cm,bottom=1cm,left=1cm,right=1cm,headheight=14pt,footskip=0.6cm}

\RequirePackage{pgfplots}
\pgfplotsset{compat=1.18}
\RequirePackage{tkz-tab}

\newcounter{partiecours}
\renewcommand{\thepartiecours}{\Roman{partiecours}}
\newcommand{\partiecours}[1]{%
  \refstepcounter{partiecours}%
  \Needspace*{8\baselineskip}%
  \bigskip
  \noindent\begin{tcolorbox}[enhanced,colback=coulPart,colframe=coulPart,
      boxrule=0pt,arc=2pt,left=8pt,right=8pt,top=4pt,bottom=4pt,
      before skip=6pt,after skip=10pt]
    \color{white}\bfseries\large \thepartiecours{} --- #1%
    \signet[partie-\arabic{partiecours}]{1}{\thepartiecours{} --- #1}
  \end{tcolorbox}%
  \addcontentsline{toc}{section}{\thepartiecours{} --- #1}%
}

\newtcolorbox{boitecours}[3][]{%
  enhanced,breakable,
  colback=white,colframe=#2,coltitle=white,
  fonttitle=\bfseries,
  attach boxed title to top left={xshift=8pt,yshift=-\tcboxedtitleheight/2},
  boxed title style={colback=#2,arc=2pt,boxrule=0pt},
  boxrule=0.9pt,arc=2pt,
  left=8pt,right=8pt,top=8pt,bottom=6pt,
  title={#3},#1}

\newcounter{methode}
\newenvironment{definitionbox}[1][Définition]%
  {\begin{boitecours}[phantom={\signet{2}{#1}}]{coulDef}{#1}}{\end{boitecours}}
\newenvironment{proprietebox}[1][Propriété]%
  {\begin{boitecours}[phantom={\signet{2}{#1}}]{coulProp}{#1}}{\end{boitecours}}
\newenvironment{methodebox}[1][Méthode]{\stepcounter{methode}%
  \begin{boitecours}[phantom={\signet[methode-\arabic{methode}]{2}{#1}}]{coulMeth}{#1}}%
  {\end{boitecours}}

\newcommand{\etape}[1]{\par\smallskip\textbf{\color{coulMeth}#1}\par\nobreak\smallskip}
\newcommand{\enonce}[1]{\emph{#1}\par\smallskip}
\newcommand{\reponse}[1]{\par\smallskip
  {\footnotesize\itshape\color{black!55}Réponses : #1}}
\newcommand{\demo}[1]{\textbf{\color{coulMeth}Démonstration.} #1}
\newcommand{\afaire}[1]{%
  \par\smallskip
  \noindent{\small\color{coulExo}$\blacktriangleright$\ \textbf{À faire :}
    fiche d'exercices du chapitre, #1.}\par\smallskip}

\newenvironment{calculs}{\setlength{\jot}{5pt}\@fleqntrue\@mathmargin\retraitcalculs
  \setlength{\abovedisplayskip}{4pt}\setlength{\belowdisplayskip}{4pt}%
  \setlength{\abovedisplayshortskip}{4pt}\setlength{\belowdisplayshortskip}{4pt}%
  \csname alignat*\endcsname{2}}%
  {\csname endalignat*\endcsname}
\newcommand{\rem}[1]{\text{\small\itshape\color{black!60}#1}}
\newcommand{\explique}[2]{\par\smallskip\noindent
  \begin{minipage}[t]{0.57\linewidth}\raggedright #1\end{minipage}\hfill
  \begin{minipage}[t]{0.40\linewidth}\raggedright\leavevmode
    {\small\itshape\color{black!60}#2}\end{minipage}\par}
\newcommand{\cas}[1]{\par\smallskip\noindent\textbf{\color{coulExo}#1}\nobreak}
\newcommand{\calc}[1]{\par\smallskip\textbf{\color{coulExo}#1}\par\nobreak}

\newtcolorbox{remarquebox}[1][Remarque]{%
  enhanced,breakable,colback=white,colframe=coulRem,
  boxrule=0pt,leftrule=2.5pt,arc=0pt,outer arc=0pt,
  left=8pt,right=6pt,top=5pt,bottom=5pt,
  before upper={\textbf{\color{coulRem}#1}\par\smallskip}}

\newtcolorbox{exemplebox}[1][Exemple]{%
  enhanced,breakable,colback=white,colframe=coulEx,
  boxrule=0pt,leftrule=2.5pt,arc=0pt,outer arc=0pt,
  left=8pt,right=6pt,top=5pt,bottom=5pt,
  before upper={\textbf{\color{coulEx}#1}\par\smallskip}}

\pgfplotsset{
  repere/.style={
    axis lines=middle,
    axis line style={-{Stealth[length=2.4mm]},thick},
    every axis x label/.style={at={(ticklabel* cs:1.0)},anchor=west},
    every axis y label/.style={at={(ticklabel* cs:1.0)},anchor=south},
    label style={font=\footnotesize},
    tick label style={font=\scriptsize},
    grid=both,
    major grid style={grille},
    minor tick num=0,
    tick align=inside,
    clip mode=individual,
  },
}
\tikzset{
  courbe/.style={coulCourbe,line width=1.1pt,smooth},
  courbeB/.style={coulCourbeB,line width=1.1pt,smooth},
  courbeC/.style={coulCourbeC,line width=1.1pt,smooth},
}

\newlist{etapes}{enumerate}{1}
\setlist[etapes]{label=\textbf{\arabic*.},leftmargin=*,itemsep=1pt,topsep=2pt}

\newlist{expressions}{enumerate}{1}
\setlist[expressions]{label={\Alph*\ $=$},leftmargin=*,itemsep=3pt,topsep=3pt}

\newcounter{exo}
\renewcommand{\theexo}{\ifnum\value{exo}<10 0\fi\arabic{exo}}
\newtcolorbox{exercice}[1][]{%
  enhanced,breakable,
  colback=white,colframe=coulExo,
  boxrule=0.9pt,arc=2pt,
  left=8pt,right=8pt,top=8pt,bottom=6pt,
  coltitle=white,fonttitle=\bfseries,
  attach boxed title to top left={xshift=8pt,yshift=-\tcboxedtitleheight/2},
  boxed title style={colback=coulExo,arc=2pt,boxrule=0pt},
  phantom={\signet{2}{Exercice \arabic{exo}}},
  title={Exercice \theexo\ifx\relax#1\relax\else{} --- #1\fi}}
\newenvironment{exo}[1][\relax]{\refstepcounter{exo}\begin{exercice}[#1]}{\end{exercice}}

  \newenvironment{correction}{\par}{\par}
  \newcommand{\autableau}{}
  \newcommand{\etapeexemples}{\etape{Exemples corrigés}}
  \newcommand{\etapetableau}[1]{\etape{#1}}
  \newenvironment{exempletableau}[1][Exemple]%
    {\begin{exemplebox}[#1]}{\end{exemplebox}}

\RequirePackage{comment}
\excludecomment{correctionavj}

\newcommand{\coursEntete}{}
\newcommand{\coursNiveau}{}
\newcommand{\coursTitre}{}
\newcommand{\coursSurTitre}{}
\newcommand{\coursSousTitre}{}

\newcommand{\enteteCours}[1]{\renewcommand{\coursEntete}{#1}}
\newcommand{\chapitrecours}[2]{\enteteCours{Chapitre #1 --- #2}}
\newcommand{\niveaucours}[1]{\renewcommand{\coursNiveau}{#1}}
\newcommand{\titrecours}[3][]{%
  \renewcommand{\coursTitre}{#2}%
  \renewcommand{\coursSurTitre}{#1}%
  \renewcommand{\coursSousTitre}{#3}}

\pagestyle{fancy}
\fancyhf{}
\fancyhead[L]{\footnotesize\color{coulPart}\textbf{\coursEntete}}
\fancyhead[R]{\footnotesize\color{coulPart}\textbf{\coursNiveau}}
\fancyfoot[C]{\footnotesize\thepage}
\renewcommand{\headrulewidth}{0.4pt}
\renewcommand{\headrule}{\hbox to\headwidth{\color{coulPart!40}\leaders\hrule height \headrulewidth\hfill}}

\setlength{\parskip}{2pt}

\AtBeginDocument{%
  \begin{center}
    \begin{tcolorbox}[enhanced,width=0.72\linewidth,colback=white,colframe=coulPart,
        boxrule=1.6pt,arc=3pt,halign=center,top=10pt,bottom=10pt]
      \ifx\coursSurTitre\empty
        {\Huge\bfseries\color{coulPart} \coursTitre}\\[4pt]
      \else
        {\Huge\bfseries\color{coulPart} \coursTitre}\\[3pt]
        {\Large\bfseries\color{coulPart} \coursSurTitre}\\[5pt]
      \fi
      {\normalsize \coursSousTitre}%
      
    \end{tcolorbox}
  \end{center}%
}
\makeatother

\chapitrecours{3}{Vecteurs}
\niveaucours{Seconde}
\titrecours{VECTEURS}{Cours --- Seconde}

\begin{document}

\ShorthandsOff
\tikzset{
  vecteur/.style={-{Stealth[length=2.6mm]},line width=1.1pt,coulCourbe},
  vecteurB/.style={-{Stealth[length=2.6mm]},line width=1.1pt,coulCourbeB},
  vecteurC/.style={-{Stealth[length=2.6mm]},line width=1.1pt,coulCourbeC},
  aide/.style={dashed,black!45,line width=0.6pt},
  fig/.style={line width=0.8pt,black!75},
}
\ShorthandsOn

\begin{remarquebox}[Comment utiliser ce cours]
Chaque partie commence par ce qu'il faut \textbf{savoir} (définitions, propriétés), puis
vient ce qu'il faut \textbf{savoir faire} : une méthode, avec des
\textbf{exemples corrigés} faits en classe, puis des exercices
\og\textbf{À vous de jouer}\fg{} à chercher seul. Les figures se font au crayon et à la
règle : un vecteur mal tracé est un vecteur faux.

\end{remarquebox}

\partiecours{Caractéristiques d'un vecteur}

\begin{definitionbox}[Définition 1 --- Vecteur]
\begin{minipage}[t]{0.58\linewidth}\raggedright
Deux points $A$ et $B$ définissent le \textbf{vecteur} $\vect{AB}$ : il va de son
\textbf{origine} $A$ à son \textbf{extrémité} $B$. On le représente par une flèche
de $A$ vers $B$. Il a trois caractéristiques :
\begin{itemize}
  \item une \textbf{direction}, celle de la droite $(AB)$ ;
  \item un \textbf{sens}, de $A$ vers $B$ ;
  \item une \textbf{longueur}, appelée \textbf{norme} et notée $\norme{\vect{AB}}=AB$.
\end{itemize}
On nomme aussi un vecteur par une seule lettre : $\vec{u}$, $\vec{v}$\dots
\end{minipage}\hfill
\begin{minipage}[t]{0.38\linewidth}
\vspace{0pt}\centering
\begin{tikzpicture}[scale=0.8]
  \draw[aide] (-1,-0.5) -- (5,2.5) node[right,font=\footnotesize,black]{$(AB)$};
  \draw[vecteur] (0,0) -- node[midway,above left,font=\footnotesize]{$\vect{AB}$} (4,2);
  \pt{0,0}{below right}{A} \pt{4,2}{below right}{B}
  \node[font=\footnotesize,below=0.35cm] at (0,0) {origine};
  \node[font=\footnotesize,below=0.35cm] at (4,2) {extrémité};
\end{tikzpicture}
\end{minipage}
\end{definitionbox}

\begin{remarquebox}[Deux confusions à éviter]
\begin{itemize}
  \item La \textbf{direction} n'est pas le \textbf{sens} : une direction a deux sens.
        $\vect{AB}$ et $\vect{BA}$ ont la même direction et des sens contraires.
  \item $AB$ est un nombre (une distance), $\vect{AB}$ est un vecteur : on écrit
        $\norme{\vect{AB}}=3$, jamais $\vect{AB}=3$.
\end{itemize}
\end{remarquebox}

\begin{definitionbox}[Définition 2 --- Vecteurs égaux, nul, opposés]
\begin{itemize}
  \item Deux vecteurs sont \textbf{égaux} s'ils ont même direction, même sens et même
        norme. Un vecteur peut donc être copié et déplacé : on dit qu'il a une infinité
        de \textbf{représentants}.
  \item Le vecteur $\vect{AA}$ est le \textbf{vecteur nul}, noté $\vec{0}$.
  \item Deux vecteurs sont \textbf{opposés} s'ils ont même direction, même norme et
        des sens contraires : $\vect{BA}=-\vect{AB}$.
\end{itemize}
\end{definitionbox}

\begin{exemplebox}[Exemple --- Trois représentants d'un même vecteur]
\begin{center}
\begin{tikzpicture}[scale=0.55]
  \quadrillage{-0.4}{-0.4}{11.4}{3.4}
  \draw[vecteur] (0,0) -- (3,1);  \node[font=\small,coulCourbe] at (1.2,1) {$\vec{u}$};
  \draw[vecteur] (4,2) -- (7,3);  \node[font=\small,coulCourbe] at (5.2,3) {$\vec{u}$};
  \draw[vecteur] (8,0) -- (11,1); \node[font=\small,coulCourbe] at (9.2,1) {$\vec{u}$};
\end{tikzpicture}
\end{center}
Les trois flèches font le même déplacement ($3$ carreaux à droite et $1$ vers le haut) :
ce sont trois représentants du même vecteur $\vec{u}$.
\end{exemplebox}

\begin{methodebox}[Méthode 1 --- Reconnaître des vecteurs égaux ou opposés]
Sur un quadrillage, on lit le \textbf{déplacement} de chaque flèche : tant de carreaux
horizontalement, tant de carreaux verticalement.
\etapeexemples
\begin{questions}
  \item Tracer les vecteurs $\vect{AB}$, $\vect{CD}$, $\vect{EF}$ et $\vect{GH}$.
  \item Pour chaque vecteur $\vect{CD}$, $\vect{EF}$ et $\vect{GH}$, dire s'il est égal à
        $\vect{AB}$, opposé à $\vect{AB}$, ou ni l'un ni l'autre.
\end{questions}
\begin{center}
\begin{tikzpicture}[scale=0.6]
  \quadrillage{-0.4}{-0.4}{10.4}{3.4}
  \pt{0,0}{below left}{A} \pt{4,1}{below right}{B}
  \pt{1,2}{above left}{C} \pt{5,3}{above right}{D}
  \pt{10,3}{above right}{E} \pt{6,2}{above left}{F}
  \pt{7,0}{below left}{G} \pt{10,1}{below right}{H}
\end{tikzpicture}
\end{center}
\begin{correction}
\cas{1\textdegree)}\par\nopagebreak
\begin{center}
\begin{tikzpicture}[scale=0.6]
  \quadrillage{-0.4}{-0.4}{10.4}{3.4}
  \draw[vecteur] (0,0) -- node[midway,above,font=\footnotesize]{$\vect{AB}$} (4,1);
  \draw[vecteur] (1,2) -- node[midway,above,font=\footnotesize]{$\vect{CD}$} (5,3);
  \draw[vecteurB] (10,3) -- node[midway,above,font=\footnotesize]{$\vect{EF}$} (6,2);
  \draw[vecteurC] (7,0) -- node[midway,above,font=\footnotesize]{$\vect{GH}$} (10,1);
  \pt{0,0}{below left}{A} \pt{4,1}{below right}{B}
  \pt{1,2}{above left}{C} \pt{5,3}{above right}{D}
  \pt{10,3}{above right}{E} \pt{6,2}{above left}{F}
  \pt{7,0}{below left}{G} \pt{10,1}{below right}{H}
\end{tikzpicture}
\end{center}
\explique{Pour $\vect{AB}$, on trace une flèche qui part de $A$ et dont la pointe est en
$B$.}{l'origine est la première lettre}
\explique{De même, $\vect{CD}$ va de $C$ à $D$, $\vect{EF}$ de $E$ à $F$, $\vect{GH}$ de
$G$ à $H$.}{la pointe est sur la seconde lettre}
\cas{2\textdegree)}\par\nopagebreak
\explique{$\vect{AB}$ : $4$ carreaux à droite et $1$ vers le haut.}{le vecteur de
référence}
\explique{$\vect{CD}$ : $4$ à droite, $1$ vers le haut, donc $\vect{CD}=\vect{AB}$.}{même
déplacement}
\explique{$\vect{EF}$ : $4$ à gauche, $1$ vers le bas, donc $\vect{EF}=-\vect{AB}$.}{déplacement
contraire}
\explique{$\vect{GH}$ : $3$ à droite, $1$ vers le haut : ni égal, ni opposé.}{autre
déplacement}
\end{correction}

\etape{À vous de jouer}
\begin{center}
\begin{tikzpicture}[scale=0.6]
  \quadrillage{-0.4}{-0.4}{11.4}{4.4}
  \pt{0,0}{below left}{M} \pt{3,2}{above left}{N}
  \pt{4,0}{below right}{P} \pt{7,2}{right}{Q}
  \pt{7,4}{above right}{R} \pt{4,2}{above left}{S}
  \pt{8,0}{below left}{T} \pt{11,1}{below right}{U}
\end{tikzpicture}
\end{center}
\begin{questions}
  \item Tracer les vecteurs $\vect{MN}$, $\vect{PQ}$, $\vect{RS}$ et $\vect{TU}$.
  \item Pour chaque vecteur $\vect{PQ}$, $\vect{RS}$ et $\vect{TU}$, dire s'il est égal à
        $\vect{MN}$, opposé à $\vect{MN}$, ou ni l'un ni l'autre.
  \item Comparer $\vect{SR}$ et $\vect{MN}$.
  \item Tracer le représentant de $\vect{MN}$ d'origine $T$.
\end{questions}
\reponse{2\textdegree) $\vect{PQ}=\vect{MN}$, $\vect{RS}=-\vect{MN}$, $\vect{TU}$ ni l'un
ni l'autre.\; 3\textdegree) $\vect{SR}=\vect{MN}$ : échanger les lettres change le sens.\;
4\textdegree) depuis $T$, $3$ carreaux à droite et $2$ vers le haut.}
\end{methodebox}

\afaire{exercices 1 à 5}

\partiecours{Translation}

\begin{definitionbox}[Définition 3 --- Translation de vecteur $\vec{u}$]
La \textbf{translation de vecteur $\vec{u}$}, notée $t_{\vec{u}}$, fait glisser chaque
point $M$ jusqu'au point $M'$ tel que $\vect{MM'}=\vec{u}$. Le point $M'$ est
l'\textbf{image} de $M$ par $t_{\vec{u}}$.
\end{definitionbox}

\begin{proprietebox}[Propriété 1]
Une translation conserve les longueurs, les angles et les aires : une figure et son
image sont superposables. Une droite a pour image une droite \textbf{parallèle}.
\end{proprietebox}

\begin{methodebox}[Méthode 2 --- Construire l'image d'un point, d'une figure]
On reporte le déplacement de $\vec{u}$ à partir du point à transformer ; pour une
figure, on transforme chaque sommet, puis on relie les images dans le même ordre.
\etapeexemples
Construire l'image $A'B'C'$ du triangle $ABC$ par la translation de vecteur $\vec{u}$.
\begin{center}
\begin{tikzpicture}[scale=0.6]
  \quadrillage{-0.4}{-0.4}{9.4}{5.4}
  \draw[vecteur] (4,5) -- (9,4);
  \node[font=\small,coulCourbe] at (6.5,4.85) {$\vec{u}$};
  \draw[fig] (0,2) -- (2,3) -- (1,4) -- cycle;
  \pt{0,2}{below left}{A} \pt{2,3}{right}{B} \pt{1,4}{above left}{C}

  \draw[fig,coulCourbeC] (5,1) -- (7,2) -- (6,3) -- cycle;
  \pt{5,1}{below left}{A'} \pt{7,2}{right}{B'} \pt{6,3}{above left}{C'}
  \draw[vecteur,line width=0.8pt] (0,2) -- node[midway,below,font=\footnotesize]{$\vect{AA'}$} (5,1);
  \draw[vecteur,line width=0.8pt] (2,3) -- node[midway,below,font=\footnotesize]{$\vect{BB'}$} (7,2);
  \draw[vecteur,line width=0.8pt] (1,4) -- node[midway,above,font=\footnotesize]{$\vect{CC'}$} (6,3);

\end{tikzpicture}
\end{center}
\begin{correction}
\cas{Résolution}\par\nopagebreak
\explique{$\vec{u}$ : $5$ carreaux à droite et $1$ vers le bas.}{on lit le déplacement}
\explique{On le reporte depuis $A$, depuis $B$ et depuis $C$ : $\vect{AA'}$,
$\vect{BB'}$ et $\vect{CC'}$ sont égaux à $\vec{u}$.}{chaque sommet glisse de la même façon}
\end{correction}

\etape{À vous de jouer}
Construire l'image $D'E'F'$ du triangle $DEF$ par la translation de vecteur $\vec{v}$.
\begin{center}
\begin{tikzpicture}[scale=0.6]
  \quadrillage{-0.4}{-0.4}{12.4}{5.4}
  \draw[vecteur] (8,0) -- node[midway,below right,font=\footnotesize]{$\vec{v}$} (12,2);
  \draw[fig] (0,1) -- (3,0) -- (2,3) -- cycle;
  \pt{0,1}{left}{D} \pt{3,0}{below}{E} \pt{2,3}{above left}{F}
\end{tikzpicture}
\end{center}
\reponse{$\vec{v}$ : $4$ carreaux à droite et $2$ vers le haut ; on le reporte depuis
$D$, $E$ et $F$.}
\end{methodebox}

\afaire{exercices 6 à 10}

\partiecours{Parallélogramme et égalité de vecteurs}

\begin{proprietebox}[Théorème 1 --- Parallélogramme]
Pour quatre points $A$, $B$, $C$, $D$ : \quad
$\vect{AB}=\vect{CD}$ \quad équivaut à \quad $ABDC$ est un parallélogramme
(éventuellement aplati, si les points sont alignés).
\begin{center}
\begin{tikzpicture}[scale=0.85]
  \draw[fig] (0,0) -- (3,0.6) -- (4,2.4) -- (1,1.8) -- cycle;
  \draw[vecteur] (0,0) -- node[midway,below,font=\footnotesize]{$\vect{AB}$} (3,0.6);
  \draw[vecteur] (1,1.8) -- node[midway,above,font=\footnotesize]{$\vect{CD}$} (4,2.4);
  \pt{0,0}{below left}{A} \pt{3,0.6}{below right}{B}
  \pt{1,1.8}{above left}{C} \pt{4,2.4}{above right}{D}
\end{tikzpicture}
\end{center}
\end{proprietebox}

\begin{remarquebox}[L'ordre des lettres]
On nomme un quadrilatère en faisant le tour de la figure dans le sens choisi.
$\vect{AB}=\vect{CD}$ donne ainsi le parallélogramme $ABDC$, et non $ABCD$ : on va de $A$
à $B$, puis à $D$, puis à $C$.
\end{remarquebox}

\begin{methodebox}[Méthode 3 --- Placer un point défini par une égalité de vecteurs]
On lit le déplacement du vecteur connu, puis on le reporte depuis l'origine imposée par
l'égalité.
\etapeexemples
Placer le point $F$ tel que $\vect{CF}=\vect{AB}$, puis nommer le parallélogramme obtenu.
\begin{center}
\begin{tikzpicture}[scale=0.6]
  \quadrillage{-0.4}{-0.4}{8.4}{3.4}
  \draw[vecteur] (0,0) -- node[midway,above left,font=\footnotesize]{$\vect{AB}$} (3,2);
  \pt{0,0}{below left}{A} \pt{3,2}{above}{B} \pt{4,0}{below}{C}

  \draw[vecteurB] (4,0) -- node[midway,below right,font=\footnotesize]{$\vect{CF}$} (7,2);
  \pt{7,2}{above right}{F}

\end{tikzpicture}
\end{center}
\begin{correction}
\cas{Résolution}\par\nopagebreak
\explique{$\vect{AB}$ : $3$ carreaux à droite et $2$ vers le haut.}{on lit le
déplacement}
\explique{Depuis $C$, on avance de $3$ à droite et $2$ vers le haut : c'est $F$.}{on le
reporte depuis $C$}
\explique{$\vect{AB}=\vect{CF}$, donc $ABFC$ est un parallélogramme.}{théorème 1}
\end{correction}

\etape{À vous de jouer}
$ABCD$ est un parallélogramme. Construire les points $E$, $F$, $G$ et $H$ tels que
\[
  \vect{BE}=\vect{AB} \qquad \vect{CF}=\vect{BC} \qquad
  \vect{DG}=\vect{CD} \qquad \vect{HA}=\vect{AB}
\]
\reponse{$B$ est le milieu de $[AE]$ ;\; $C$ celui de $[BF]$ ;\; $D$ celui de $[CG]$ ;\;
$A$ celui de $[HB]$.}
\end{methodebox}

\afaire{exercices 11 à 14}

\partiecours{Somme de deux vecteurs}

\begin{definitionbox}[Définition 4 --- Somme]
Enchaîner la translation de vecteur $\vec{u}$ puis celle de vecteur $\vec{v}$ revient à
faire une seule translation, de vecteur $\vec{u}+\vec{v}$. On trace les deux vecteurs
\textbf{bout à bout} (l'origine du second sur l'extrémité du premier) : la somme va du
départ à l'arrivée.

\smallskip
La \textbf{différence} est la somme avec l'opposé : $\vec{u}-\vec{v}=\vec{u}+(-\vec{v})$.
\end{definitionbox}

\begin{methodebox}[Méthode 4 --- Construire une somme ou une différence]
On met les vecteurs \textbf{l'un à la suite de l'autre} : l'origine du second sur
l'extrémité du premier.
\etapeexemples
\begin{questions}[label=\alph*)]
  \item Construire $\vec{u}+\vec{v}$, puis $\vec{u}-\vec{v}$.
  \item Construire le point $D$ tel que $\vect{AD}=\vect{AB}+\vect{AC}$. Quelle est la
        nature de $ABDC$ ?
\end{questions}
\begin{center}
\begin{tikzpicture}[scale=0.55]
  \node[font=\small\bfseries,coulExo] at (3,-1.1) {a)};
  \quadrillage{-0.4}{-0.4}{6.4}{3.4}
  \draw[vecteur] (0,1) -- (3,2);  \node[font=\small,coulCourbe] at (1.3,1.9) {$\vec{u}$};
  \draw[vecteurB] (5,3) -- (6,1); \node[font=\small,coulCourbeB] at (5.9,2.3) {$\vec{v}$};
\end{tikzpicture}
\qquad\qquad
\begin{tikzpicture}[scale=0.55]
  \node[font=\small\bfseries,coulExo] at (2.5,-1.1) {b)};
  \quadrillage{-0.4}{-0.4}{5.4}{4.4}
  \draw[vecteur] (0,0) -- node[midway,below right,font=\footnotesize]{$\vect{AB}$} (4,1);
  \draw[vecteurB] (0,0) -- node[midway,left,font=\footnotesize]{$\vect{AC}$} (1,3);
  \pt{0,0}{below left}{A} \pt{4,1}{below right}{B} \pt{1,3}{above left}{C}
\end{tikzpicture}
\end{center}
\begin{correction}
\cas{a)}\par\nopagebreak
\begin{center}
\begin{tikzpicture}[scale=0.6]
  \quadrillage{-0.4}{-2.4}{10.4}{3.4}
  \draw[vecteur] (0,1) -- (3,2);
  \draw[vecteurB] (3,2) -- (4,0);
  \draw[vecteurC] (0,1) -- (4,0);
  \node[font=\small,coulCourbe] at (1.3,1.9) {$\vec{u}$};
  \node[font=\small,coulCourbeB] at (3.9,1.2) {$\vec{v}$};
  \node[font=\small,coulCourbeC] at (1.8,0.1) {$\vec{u}+\vec{v}$};
  \draw[vecteur] (6,0) -- (9,1);
  \draw[vecteurB] (9,1) -- (8,3);
  \draw[vecteurC] (6,0) -- (8,3);
  \node[font=\small,coulCourbe] at (7.6,0.1) {$\vec{u}$};
  \node[font=\small,coulCourbeB] at (9.2,2.2) {$-\vec{v}$};
  \node[font=\small,coulCourbeC] at (6.3,1.9) {$\vec{u}-\vec{v}$};
\end{tikzpicture}
\end{center}
\explique{On lit les déplacements : $\vec{u}$, $3$ carreaux à droite et $1$ vers le haut ;
$\vec{v}$, $1$ à droite et $2$ vers le bas.}{on compte les carreaux}
\explique{On trace $\vec{u}$, puis $\vec{v}$ à la suite : $\vec{u}+\vec{v}$ fait $4$
carreaux à droite et $1$ vers le bas.}{l'un à la suite de l'autre}
\explique{On trace $\vec{u}$, puis $-\vec{v}$ ($1$ à gauche, $2$ vers le haut) :
$\vec{u}-\vec{v}$ fait $2$ carreaux à droite et $3$ vers le haut.}{on ajoute l'opposé}
\cas{b)}\par\nopagebreak
\begin{center}
\begin{tikzpicture}[scale=0.55]
  \quadrillage{-0.4}{-0.4}{5.4}{4.4}
  \draw[fig,black!30] (1,3) -- (5,4);
  \draw[vecteur] (0,0) -- node[midway,below right,font=\footnotesize]{$\vect{AB}$} (4,1);
  \draw[vecteurB] (0,0) -- node[midway,left,font=\footnotesize]{$\vect{AC}$} (1,3);
  \draw[vecteurB] (4,1) -- node[midway,right,font=\footnotesize]{$\vect{BD}$} (5,4);
  \draw[vecteurC] (0,0) -- node[midway,above left,font=\footnotesize]{$\vect{AD}$} (5,4);
  \pt{0,0}{below left}{A} \pt{4,1}{below right}{B} \pt{1,3}{above left}{C}
  \pt{5,4}{above right}{D}
\end{tikzpicture}
\end{center}
\explique{$\vect{AC}$ : $1$ carreau à droite et $3$ vers le haut ; on le reporte à la
suite de $\vect{AB}$, depuis $B$ : on arrive en $D$.}{l'un à la suite de l'autre}
\explique{$\vect{BD}=\vect{AC}$, donc $ABDC$ est un parallélogramme.}{théorème 1}
\explique{Quand deux vecteurs ont la même origine, leur somme suit la diagonale du
parallélogramme.}{à retenir}
\end{correction}

\etape{À vous de jouer}
\begin{minipage}[t]{0.56\linewidth}\raggedright
Reproduire le triangle $ABC$ ci-contre.
\begin{questions}
  \item Construire le point $D$ tel que $\vect{BD}=\vect{BA}+\vect{BC}$. Quelle est la
        nature de $BADC$ ?
  \item Construire le point $E$ tel que $\vect{BE}=\vect{BA}-\vect{BC}$.
\end{questions}
\end{minipage}\hfill
\begin{minipage}[t]{0.40\linewidth}
\vspace{0pt}\centering
\begin{tikzpicture}[scale=0.42]
  \quadrillage{-3.4}{-3.4}{5.4}{4.4}
  \draw[fig] (0,0) -- (5,1) -- (2,4) -- cycle;
  \pt{0,0}{below left}{A} \pt{5,1}{right}{B} \pt{2,4}{above}{C}
\end{tikzpicture}
\end{minipage}
\reponse{1\textdegree) depuis $B$, $8$ carreaux à gauche et $2$ vers le haut ; $BADC$ est
un parallélogramme.\; 2\textdegree) $\vect{BA}-\vect{BC}=\vect{BA}+\vect{CB}=\vect{CA}$ :
depuis $B$, $2$ carreaux à gauche et $4$ vers le bas.}
\end{methodebox}

\afaire{exercices 15 à 20}

\partiecours{Milieu d'un segment}

\begin{proprietebox}[Théorème 2 --- Milieu]
$I$ est le milieu de $[AB]$ \quad équivaut à \quad $\vect{AI}=\vect{IB}$.
\begin{center}
\begin{tikzpicture}[scale=0.85]
  \draw[vecteur] (0,0) -- node[midway,above left,font=\footnotesize]{$\vect{AI}$} (2.5,0.8);
  \draw[vecteur] (2.5,0.8) -- node[midway,above left,font=\footnotesize]{$\vect{IB}$} (5,1.6);
  \pt{0,0}{below left}{A} \pt{2.5,0.8}{below right}{I} \pt{5,1.6}{below right}{B}
\end{tikzpicture}
\end{center}
Pour démontrer que $M$ est le milieu de $[XY]$, il suffit donc d'obtenir
$\vect{XM}=\vect{MY}$.

\smallskip
On l'écrit aussi $\vect{IA}=-\vect{IB}$, ou $\vect{IA}+\vect{IB}=\vec{0}$. Rappel : les
diagonales d'un parallélogramme se coupent en leur milieu.
\end{proprietebox}

\begin{methodebox}[Méthode 5 --- Démontrer qu'un point est un milieu]
On traduit chaque parallélogramme de l'énoncé par une égalité de vecteurs, puis on les
enchaîne jusqu'à obtenir $\vect{XM}=\vect{MY}$.
\etapeexemples
$ABCD$ et $BCFD$ sont deux parallélogrammes. Démontrer que $D$ est le milieu de $[AF]$.
\begin{center}
\begin{tikzpicture}[scale=0.85]
  \coordinate (A) at (0,0);
  \coordinate (B) at (3,0);
  \coordinate (C) at (4,1.6);
  \coordinate (D) at (1,1.6);
  \coordinate (F) at (2,3.2);
  \draw[fig] (A) -- (B) -- (C) -- (D) -- cycle;
  \draw[fig,black!45] (B) -- (C) -- (F) -- (D) -- cycle;
  \pt{A}{below left}{A} \pt{B}{below right}{B} \pt{C}{right}{C}
  \pt{D}{left}{D} \pt{F}{above}{F}
\end{tikzpicture}
\end{center}
\begin{correction}
\cas{Résolution}\par\nopagebreak
\explique{$ABCD$ est un parallélogramme, donc $\vect{AD}=\vect{BC}$.}{côtés opposés}
\explique{$BCFD$ est un parallélogramme, donc $\vect{BC}=\vect{DF}$.}{côtés opposés}
\explique{Ainsi $\vect{AD}=\vect{DF}$ : $D$ est le milieu de $[AF]$.}{théorème 2}
\end{correction}

\etape{À vous de jouer}
$ABCD$ et $ABEC$ sont deux parallélogrammes. Démontrer que $C$ est le milieu de $[DE]$.
\begin{center}
\begin{tikzpicture}[scale=0.85]
  \coordinate (A) at (0,0);
  \coordinate (B) at (3,0);
  \coordinate (C) at (4,1.6);
  \coordinate (D) at (1,1.6);
  \coordinate (E) at (7,1.6);
  \draw[fig] (A) -- (B) -- (C) -- (D) -- cycle;
  \draw[fig,black!45] (A) -- (B) -- (E) -- (C) -- cycle;
  \pt{A}{below left}{A} \pt{B}{below right}{B} \pt{C}{above}{C}
  \pt{D}{above left}{D} \pt{E}{above right}{E}
\end{tikzpicture}
\end{center}
\reponse{$\vect{DC}=\vect{AB}$ et $\vect{AB}=\vect{CE}$, donc $\vect{DC}=\vect{CE}$ : $C$
est le milieu de $[DE]$.}
\end{methodebox}

\afaire{exercices 21 à 23}

\partiecours{Multiplier un vecteur par un réel}

\begin{definitionbox}[Définition 5 --- Produit par un réel]
Pour un vecteur $\vec{u}$ non nul et un réel $k$ non nul, le vecteur $k\vec{u}$ :
\begin{itemize}
  \item a la \textbf{même direction} que $\vec{u}$ ;
  \item a le \textbf{même sens} si $k>0$, le \textbf{sens contraire} si $k<0$ ;
  \item est $k$ fois plus long si $k>0$ (et $-k$ fois plus long si $k<0$).
\end{itemize}
Si $k=0$ ou $\vec{u}=\vec{0}$, alors $k\vec{u}=\vec{0}$. Par exemple,
$\vec{u}+\vec{u}+\vec{u}=3\vec{u}$ et $(-1)\vec{u}=-\vec{u}$.
\end{definitionbox}

\begin{remarquebox}[À retenir]
Si $I$ est le milieu de $[AB]$ : $\vect{AI}=\frac12\vect{AB}$ et $\vect{AB}=2\vect{AI}$.
On développe comme en calcul littéral : $k(\vec{u}+\vec{v})=k\vec{u}+k\vec{v}$.
\end{remarquebox}

\begin{methodebox}[Méthode 6 --- Construire $k\vec{u}$]
On multiplie par $k$ les deux nombres de carreaux du déplacement de $\vec{u}$ ; si $k<0$,
on change le sens.
\etapeexemples
$\vec{u}$ : $2$ carreaux à droite et $2$ vers le haut. Tracer $2\vec{u}$ et $-1{,}5\vec{u}$.
\begin{correction}
\begin{center}
\begin{tikzpicture}[scale=0.55]
  \quadrillage{-3.4}{-3.4}{8.4}{4.4}
  \draw[vecteur] (-3,0) -- (-1,2);
  \node[font=\small,coulCourbe] at (-2.6,1.5) {$\vec{u}$};
  \draw[vecteurB] (0,-2) -- (4,2);
  \node[font=\small,coulCourbeB] at (3.2,0.6) {$2\vec{u}$};
  \draw[vecteurC] (8,4) -- (5,1);
  \node[font=\small,coulCourbeC] at (5.4,2.8) {$-1{,}5\vec{u}$};
\end{tikzpicture}
\end{center}
\explique{$2\vec{u}$ : $4$ carreaux à droite et $4$ vers le haut.}{même sens, deux fois
plus long}
\explique{$-1{,}5\vec{u}$ : $3$ carreaux à gauche et $3$ vers le bas.}{sens contraire,
une fois et demie plus long}
\end{correction}

\etape{À vous de jouer}
$\vec{v}$ : $2$ carreaux à droite et $2$ vers le bas. Tracer $3\vec{v}$ et $-1{,}5\vec{v}$.
\reponse{$3\vec{v}$ : $6$ carreaux à droite et $6$ vers le bas ;\; $-1{,}5\vec{v}$ :
$3$ carreaux à gauche et $3$ vers le haut.}
\end{methodebox}

\afaire{exercices 24 à 26}

\partiecours{Relation de Chasles}

\begin{proprietebox}[Théorème 3 --- Relation de Chasles]
Pour tous points $A$, $B$ et $C$ : \quad $\vect{AB}+\vect{BC}=\vect{AC}$.

\begin{itemize}
  \item Aller de $A$ à $B$, puis de $B$ à $C$, c'est aller de $A$ à $C$.
  \item $\vect{AB}+\vect{BA}=\vect{AA}=\vec{0}$, donc $-\vect{AB}=\vect{BA}$.
  \item On peut passer par n'importe quel point de son choix :
        $\vect{AC}=\vect{AM}+\vect{MC}$.
\end{itemize}
\begin{center}
\begin{tikzpicture}[scale=0.85]
  \draw[vecteur] (0,0) -- node[midway,above left,font=\footnotesize]{$\vect{AB}$} (2.6,1.6);
  \draw[vecteurB] (2.6,1.6) -- node[midway,above right,font=\footnotesize]{$\vect{BC}$} (5.5,0.7);
  \draw[vecteurC] (0,0) -- node[midway,below,font=\footnotesize]{$\vect{AC}$} (5.5,0.7);
  \pt{0,0}{below left}{A} \pt{2.6,1.6}{above}{B} \pt{5.5,0.7}{below right}{C}
\end{tikzpicture}
\end{center}
\end{proprietebox}

\begin{methodebox}[Méthode 7 --- Simplifier une somme de vecteurs]
On remplace chaque $-\vect{XY}$ par $\vect{YX}$, puis on range les vecteurs pour que
l'extrémité de chacun soit l'origine du suivant.
\etapeexemples
Simplifier :
\begin{questions}[label=\alph*)]
  \item $\vect{KL}+\vect{JK}$
  \item $\vect{AB}+\vect{DA}+\vect{CD}$
  \item $\vect{PQ}-\vect{RQ}+\vect{RP}$
\end{questions}
\begin{correction}
\cas{a)}
\begin{calculs}
\vect{KL}+\vect{JK} &= \vect{JK}+\vect{KL} &\qquad& \rem{on range : $J\to K\to L$}\\
\vect{KL}+\vect{JK} &= \vect{JL} &\qquad& \rem{Chasles}
\end{calculs}
\cas{b)}
\begin{calculs}
\vect{AB}+\vect{DA}+\vect{CD} &= \vect{CD}+\vect{DA}+\vect{AB} &\qquad& \rem{on range : $C\to D\to A\to B$}\\
\vect{AB}+\vect{DA}+\vect{CD} &= \vect{CB} &&
\end{calculs}
\cas{c)}
\begin{calculs}
\vect{PQ}-\vect{RQ}+\vect{RP} &= \vect{PQ}+\vect{QR}+\vect{RP} &\qquad& \rem{$-\vect{RQ}=\vect{QR}$}\\
\vect{PQ}-\vect{RQ}+\vect{RP} &= \vect{PP}=\vec{0} &\qquad& \rem{on revient au point de départ}
\end{calculs}
\end{correction}

\begin{remarquebox}[Erreur à éviter]
Chasles ne s'applique que si la lettre d'arrivée du premier vecteur est la lettre de
départ du second : $\vect{AB}+\vect{CD}$ ne se simplifie pas.
\end{remarquebox}

\etape{À vous de jouer}
Simplifier :
\[
  \vect{UV}+\vect{VW} \qquad \vect{GH}+\vect{FG} \qquad \vect{MN}+\vect{NP}+\vect{PM}
  \qquad \vect{AC}+\vect{BA}+\vect{CB} \qquad \vect{EF}-\vect{GF}
\]
\reponse{$\vect{UW}$ ;\; $\vect{FH}$ ;\; $\vec{0}$ ;\; $\vect{BA}+\vect{AC}+\vect{CB}=\vec{0}$ ;\;
$\vect{EF}+\vect{FG}=\vect{EG}$.}
\end{methodebox}

\afaire{exercices 27 à 33}

\end{document}
